\begin{enumerate}
%1
\item \textbf{การคำนวณความชันของเส้นโค้ง}

กำหนดฟังก์ชัน $f(x) = \sin x$ จงใช้เครื่องมือตามลิงก์นี้ เพื่อตอบคำถามต่อไปนี้

\begin{center}
\href{https://www.desmos.com/calculator/zqfzwqzo4n}{https://www.desmos.com/calculator/zqfzwqzo4n} 
\end{center}
\begin{tasks}(1)
\task วาดกราฟของ $y = f(x)$ โดยคร่าว ในช่วง $-1 \leq x \leq 1$ 
\task หาค่าของ $f(0)$ และ $f(0.5)$ ด้วยการเปลี่ยนค่า $c$ 
\task ลากเส้นตรงที่เชื่อมระหว่างจุด $(0, f(0))$ และจุด $(0.5, f(0.5))$ 
\task ถ้าพิจารณาด้วยตาเปล่าแล้ว คิดว่าความชันของเส้นสัมผัสเส้นโค้งที่จุด $x = 0$ มีค่าเท่ากับเท่าไหร่
\task คาดเดาค่าของ $f\,'(0)$ ด้วยการเปลี่ยนค่า $h$ ให้เข้าใกล้ 0
\end{tasks}

\vspace{100pt}

%2
\item \textbf{การคำนวณอนุพันธ์โดยใช้นิยาม}

นิยามอนุพันธ์ของฟังก์ชันคือ
\[ \dyx = f\,'(x) = \rule{100pt}{.5pt} \]

จงหาอนุพันธ์ของฟังก์ชันต่อไปนี้ \Big(นั่นคือ หา $\dyx$\Big) โดยใช้นิยาม
\begin{tasks}[after-skip = 100pt, after-item-skip = 100pt](2)
\task $f(x) = 4$
\task $f(x) = x$
\task $f(x) = 2x+1$
\task $f(x) = x^3$ (คำใบ้: $(x+1)^3 = x^3+3x^2+3x+1$)
\end{tasks}

\newpage

\item \textbf{สูตรและสมบัติพื้นฐานของอนุพันธ์} จงหาอนุพันธ์ต่อไปนี้
\begin{tasks}[after-skip = 50pt, after-item-skip = 50pt](3)
\task $\dx 42$
\task $\dx x^4$
\task $\dx \left(\frac{1}{x^3}\right)$
\task $\dx \sqrt[4]{x}$
\task $\dx (4x^3 + 3x^4)$
\task $\dx \left( \frac{x^3}{3} + \frac{x^2}{2}+\frac{2}{x^2}\right)$
\end{tasks}

\item \textbf{สมบัติอนุพันธ์ของผลคูณฟังก์ชัน} จงหาอนุพันธ์ต่อไปนี้
\begin{tasks}[after-skip = 100pt, after-item-skip = 100pt](2)
\task $\dx x^3(x^2-1)$
\task $\dx (2x^2-1)(x^4+3)$
\task $\dx (x-1)(x^2+x+1)$
\task $\dx (x+1)\left(1-\frac{1}{x}\right)$
\end{tasks}

\item \textbf{สมบัติอนุพันธ์ของผลหารฟังก์ชัน} จงหาอนุพันธ์ต่อไปนี้
\begin{tasks}[after-skip = 100pt, after-item-skip = 100pt](2)
\task $\dx \left(\frac{3}{x^2+1}\right)$
\task $\dx \left(\frac{x^2}{x^2+1}\right)$
\task $\dx \left(\frac{4x-3}{2x+1}\right)$
\task $\dx \left(\frac{x^3-1}{x^2-1}\right)$
\end{tasks}
\end{enumerate}

\end{document}